%%%%%%%%%%%%%%%%%%%%%%%%%%%%%%%%%%%%%%%%%%%%%%%%%%%%%%%%%%%%%%
%% Rozdział 1: Wprowadzenie
%%%%%%%%%%%%%%%%%%%%%%%%%%%%%%%%%%%%%%%%%%%%%%%%%%%%%%%%%%%%%%

\section{Wprowadzenie}
\label{chap:introduction}

\subsection{Motywacja i kontekst}
\label{sec:motivation-context}

Języki migowe są pełnoprawnymi językami naturalnymi, posiadającymi własne zasoby
leksykalne, gramatykę i struktury morfosyntaktyczne, niezależne od języków
fonicznych, z którymi współistnieją. Wbrew powszechnemu przekonaniu nie są one ani
uniwersalne, ani nie stanowią manualnej transliteracji mowy: amerykański język
migowy (ASL), brytyjski język migowy (BSL), niemiecki język migowy (DGS) i polski
język migowy (PJM) nie są wzajemnie zrozumiałe, a gramatyka języka migowego często
nie przypomina gramatyki języka fonicznego używanego w jego otoczeniu. Na przykład
ASL nie stosuje angielskiego szyku wyrazów, a informacje gramatyczne -- negację,
pytania czy warunkowość -- przekazuje za pomocą wykładników niemanualnych, takich
jak mimika twarzy, pozycja głowy i kierunek spojrzenia, równolegle z manualną
artykulacją znaków. Dla społeczności Głuchych i osób słabosłyszących, dla których
język migowy jest podstawowym językiem komunikacji, ta odrębność stwarza
rzeczywistą barierę w porozumiewaniu się ze słyszącą większością i uzasadnia
rozwój systemów automatycznego tłumaczenia języka migowego.

Automatyczne tłumaczenie języka migowego nie jest pojedynczym zadaniem, lecz
zbiorem powiązanych ze sobą podproblemów. Tłumaczenie między językiem fonicznym
a językiem migowym może przebiegać w obu kierunkach i zazwyczaj jest dzielone
na etapy z wykorzystaniem pośredniej reprezentacji symbolicznej w postaci
\emph{glos}. Zapis glosowy jest pisemną transkrypcją sekwencji znaków, zwykle
przedstawianą za pomocą wyrazów języka fonicznego zapisanych wielkimi literami.
Określa on, jakie znaki są wykonywane i w jakiej kolejności, pomijając ich
realizację wizualną. Na przykład wypowiedź w DGS zapisana glosami jako
ICH OSTERN WETTER ZUFRIEDEN odpowiada naturalnemu zdaniu niemieckiemu
``Da k\"onnen wir mit unserem Osterwetter eigentlich ganz zufrieden sein'' --
co pokazuje, że zapis glosowy jest zwięzłą, stratną reprezentacją pośrednią,
a nie pełnym tłumaczeniem. Taki podział prowadzi do czterech zadań kierunkowych:
przekształcania wideo w glosy i glos w tekst po stronie rozpoznawania oraz
tekstu w glosy i glos w wideo z wykorzystaniem awatara po stronie generowania.
Zadania te oraz miejsce niniejszego projektu wśród nich omówiono bardziej
szczegółowo w rozdziale~\ref{chap:related-work}.

Ze względu na ograniczoną dostępność anotowanych danych do badań nad językami
migowymi oraz praktyczną potrzebę narzędzi do ich gromadzenia i oznaczania projekt
łączy prace badawcze z pracami programistycznymi. Część badawcza koncentruje się
na rozpoznawaniu glos na podstawie wideo oraz rozszerzeniu tego podejścia
w kierunku tłumaczenia ciągłych wypowiedzi na poziomie zdań. Część programistyczna
obejmuje budowę wspierającej infrastruktury oprogramowania -- platformy do
anotacji danych oraz edukacyjnej platformy demonstracyjnej -- niezbędnej do
gromadzenia danych treningowych i prezentowania opracowanych modeli użytkownikom
końcowym.

\subsection{Cele i zakres}
\label{sec:goals-scope}

Głównym celem projektu jest zbadanie, opracowanie i wdrożenie metod tłumaczenia
języka migowego na język pisany lub mowę, a także stworzenie towarzyszących im
systemów informatycznych, które zapewnią odtwarzalność badań i umożliwią
praktyczne wykorzystanie ich wyników. Ze względu na ograniczone zasoby danych
dostępne na rynku projekt koncentruje się zarówno na samym problemie badawczym,
jak i na stworzeniu kompleksowego systemu o wymiernej wartości praktycznej
i naukowej. W szczególności projekt zakłada realizację następujących celów:

\begin{itemize}
  \item zaprojektowanie i wdrożenie zdecentralizowanej platformy do anotacji
        danych, służącej do gromadzenia i oznaczania zbiorów nagrań wideo
        w języku migowym;
  \item opracowanie i ocena modelu rozpoznawania izolowanych znaków migowych
        na podstawie wideo (wideo--glosy), z porównaniem wielu architektur
        enkoderów czasowych w identycznych warunkach eksperymentalnych;
  \item zbadanie tłumaczenia między sekwencjami glos a tekstem w języku
        naturalnym w obu kierunkach (glosy--tekst i tekst--glosy);
  \item integrację komponentów rozpoznawania izolowanych znaków i tłumaczenia
        w całościowy potok przetwarzania wideo na tekst na poziomie zdań;
  \item dostarczenie platformy edukacyjnej prezentującej wytrenowane modele
        użytkownikom końcowym za pośrednictwem przystępnego interfejsu.
\end{itemize}

\todo{DO UZUPEŁNIENIA: Należy precyzyjnie określić granice zakresu pracy -- jakie
języki migowe i zbiory danych są objęte badaniami, co wyraźnie wykracza poza
zakres pracy (np. tłumaczenie ciągłe lub w czasie rzeczywistym, wdrożenie
obsługujące wiele osób migających, generowanie wypowiedzi z wykorzystaniem
awatara) oraz jakie ograniczenia wynikają z danych źródłowych (np. wyłącznie
izolowane, uprzednio wyodrębnione znaki).}

\subsection{Struktura pracy}
\label{sec:thesis-structure}

Dalsza część pracy ma następującą strukturę.
Rozdział~\ref{chap:related-work} zawiera przegląd dotychczasowych badań nad
rozpoznawaniem i tłumaczeniem języka migowego, obejmujący metody przekształcania
wideo w glosy, modele przetwarzania języka naturalnego (NLP) stosowane do języka
migowego oraz istniejące systemy i benchmarki.
Rozdział~\ref{chap:annotation-platform} opisuje platformę do anotacji stworzoną
w celu gromadzenia i oznaczania zbiorów nagrań wideo wykorzystywanych w pracy.
Rozdział~\ref{chap:isolated-gloss-recognition} przedstawia model rozpoznawania
izolowanych glos (wideo--glosy), w tym zbiór danych, potok przetwarzania wstępnego,
architekturę modelu i wyniki eksperymentów.
Rozdział~\ref{chap:text-translation} dotyczy tłumaczenia między sekwencjami glos
a tekstem w języku naturalnym.
Rozdział~\ref{chap:sentence-level-translation} opisuje integrację komponentów
z rozdziałów~\ref{chap:isolated-gloss-recognition} i~\ref{chap:text-translation}
w całościowy potok przetwarzania wideo na tekst na poziomie zdań.
Rozdział~\ref{chap:educational-platform} przedstawia platformę edukacyjną
prezentującą opracowane modele użytkownikom końcowym.
Na zakończenie rozdział~\ref{chap:conclusions} podsumowuje realizację wskazanych
wyżej celów, omawia ograniczenia opracowanego systemu i nakreśla kierunki
dalszych prac.